\scanpage{72}\mspage{35}
$\overline{\underline{\mathrm V}}$\quad Récapitulation

Soit $\Lambda$ un anneau comm. non nul, $M$ un $\Lambda$-module libre de rang 2. On lui associe deux ``structures \unsure{triangulées|triangulières}'' (ens. de sommets $S$, plus \unsure{une|un} \inserted{ens. de} parties \deleted{\unsure{à}} \inserted{$F$} $\unsure{\subset|\in}\mathfrak P_3(S)$ \deleted{\unsure{ayant comme}} \inserted{de} réunion dans $S$), dont la deuxième n'est intéressante que si $2.1_\Lambda\ne0$ (autrement, elle coïncide avec la première).

Les structures triangulées \deleted{sont} en question sont ``partitionnées'' (i.e. $F\subset\mathfrak P_3(S)$ est munie d'une partition, via une application sur un ens d'indices convenable $I$) --- donc la structure envisagée peut être considérée comme un ens. de structures triangulées sur le \unsure{$=$|même} ens $S$, indexé par $I$. Les deux structures envisagées dépendent fonctoriellement de $M$, et la deuxième est le quotient de la première par $\{\pm\mathrm{id}_M\}\subset\mathrm{Gl}_\Lambda(M)$.

A)
\[
 \left\{\begin{aligned}
 \widetilde S_M&=M^{*\unsure{*}}=\{x\in M\mid\exists y\in M,\ (x,y)\in\operatorname{Rep}(M)\}\\
 \widetilde F_M&=\deleted{\unsure{\mathfrak P}}\bigl\{\{x,y,z\}\in\mathfrak P_3(\widetilde S_M)\mid
 \substack{x+y+z=0\\x\wedge y\ (=y\wedge z=z\wedge x)\in(\det M)^*}\bigr\}
 \end{aligned}\right.\tag*{(1)}
\]
\[
 \overrightarrow{\unsure{\pi|\Pi}}_M=\det M^*,\quad \pi_M=(\det M^*)/\{\pm1_\Lambda\}\tag*{(2)}
\]
% archive sheet number: 71
